% !TEX root = report.tex

\section{Introduction}
\label{sec:introduction}

Contactless sensing of human activity is valuable in healthcare, security and
ambient monitoring. Commodity Wi-Fi devices can be repurposed as radar-like
sensors: as a person moves, their body induces micro-Doppler shifts on the
Wi-Fi signal that are captured in the channel-state information (CSI) estimated
by the receiver, \emph{without recording any video}. After suitable processing,
the CSI can be turned into a \emph{Doppler spectrogram}---a
time$\,\times\,$velocity image---whose appearance changes markedly with the
activity being performed, so the recognition problem can be cast as image
classification with a convolutional neural network (CNN)~\cite{sharp}.

The central difficulty is \emph{generalization}. Different people move
differently, and different rooms (and even the same room on different days)
distort the signal differently, so a model that classifies perfectly in its
training environment may collapse on a new one. This is the challenge raised in
the project proposal, which explicitly suggests a
\emph{contrastive / self-supervised} approach to obtain subject- and
environment-independent features. Prior systems close this gap mostly through
signal processing---robust phase sanitization and multi-antenna fusion---at the
cost of a heavy, hand-engineered pipeline~\cite{sharp}; we instead ask how far a
purely \emph{representation-learning} remedy can go on top of a minimal
front-end.

In this report we tackle five-class Wi-Fi HAR under a strict cross-domain
protocol (train on a single environment, test on six unseen ones) and study
whether Supervised Contrastive Learning (SupCon)~\cite{supcon} can \emph{learn}
environment invariance rather than engineering it. The approach is
data-efficient---it reuses the same small training set and a lightweight
encoder---and its value is that the learned features can be reused by any
downstream classifier and combined with the orthogonal signal-processing tools
of the reference system. The report is organized around two clearly separated
contributions:
\begin{itemize}
  \item \textbf{Baseline (M.~Carraro).} A faithful reimplementation, in plain
  PyTorch, of the dataset pipeline and of the baseline Inception-style CNN, with
  a cross-domain evaluation protocol (train on $S1$, test on $S2$--$S7$) and a
  per-measurement-set error breakdown.
  \item \textbf{Contrastive extension (F.~Vezzani).} A Supervised Contrastive
  Learning stage (notebook~\S9): a deeper batch-normalized encoder is pre-trained
  with a supervised-contrastive loss and domain-randomization augmentations to
  learn environment-invariant features, then evaluated with a linear probe under
  the same protocol. It improves cross-domain accuracy over the baseline on
  every set.
\end{itemize}
We use the full SHARP system~\cite{sharp} as the reference upper bound and
discuss, quantitatively, why both of our models sit below it and what the
contrastive stage does and does not recover.

The remainder of the report is structured as follows.
\secref{sec:related_work} reviews the related work on Wi-Fi sensing and
contrastive representation learning. \secref{sec:processing_architecture} gives
a high-level view of the processing pipeline; \secref{sec:model} details the
signals and features, and \secref{sec:learning_framework} describes the two
learning frameworks (baseline CNN and SupCon extension). The experimental
results and the comparison against the SHARP paper are reported in
\secref{sec:results}, and \secref{sec:conclusions} draws the concluding remarks.
